\section{Testbed Specification}
\label{app:system}

The report's claims rest on two engines and the check between them. This
appendix specifies both engines, the parity harness, the cost models and the
learning interface in enough detail to rebuild them.

\subsection{The exact replay engine}

The exact engine consumes the recorded tape at its own event granularity and
asks the policy for a decision every ten minutes. Each decision is executed at
the price observed \emph{before} the interval advances, matching the latency a
live LP faces and keeping the simulator and the backtester on the same
convention. The engine's state is the pool's integer tick lattice: liquidity
is minted and burned on the exact ticks the tape reports, and a swap is
apportioned across every tick segment it crosses. Fee accrual comes from the
protocol's own fee-growth accumulators, which are 128-bit fixed-point integers
that increase monotonically with every swap; the position's fee over a step is
its liquidity multiplied by the change in the inside fee-growth value divided
by the protocol's $2^{128}$ normalisation, exactly
as the contract computes it. The ledger therefore satisfies an accounting
identity by construction,

\begin{equation}
\text{fees} + \mathrm{IL} - \text{gas} - \text{slippage}
= \text{strat\-egy excess over the opening basket},
\label{eq:app-identity}
\end{equation}
and the identity is checked at every step of every replay.

Three numerical details matter for exactness. First, liquidity arithmetic is
done in Python integers, not floats, because a large gross liquidity in
floating point carries a residue that can round a small delta to zero.
Second, the initial position is deployed before the first decision for every
policy, so the HODL benchmark has a definite token mix; the pure basket holds
the same mix as tokens. Third, a rebalance that the position cannot afford is
recorded as rejected and the position is held, mirroring an on-chain revert
rather than forcing a sale.

\subsection{The fast simulator}

Training needs roughly $10^6$ decisions per seed, which the exact engine
cannot supply within the study's budget. The simulator mirrors the same
economics in floating point under a marginal-agent assumption: the agent's
position never moves the pool price, and its fee share is its fraction of the
active liquidity while in range. This removes the segment-level replay and
raises throughput by more than an order of magnitude. The simulator runs either
on the recorded events, the mode used for every result in this report, or on a
regime-switching jump process calibrated to the training years; the calibrated
mode is not exercised here because the study keeps the evaluation on recorded
states.

\subsection{The parity harness}

Two engines are only useful together if their disagreement is bounded. The
harness feeds both the same prices, positions and gas series and compares fees,
impermanent loss, PnL, gas and slippage per step against a tolerance of
$10^{-3}$ USDC. On the recorded data the residual fee drift is on the order of
$10^{-7}$ USDC per step per seed, recorded with the engine validation
($8.9\times10^{-8}$, $1.4\times10^{-7}$, $1.0\times10^{-7}$,
$1.1\times10^{-7}$ for four seeds); impermanent loss matches exactly, and the
pooled error is dominated by a single seed's tail but remains two orders of
magnitude inside tolerance.

A separate suite of thirteen look-ahead probes attacks the composed system at
the seams where a look-ahead error would be easiest to hide: the split wall for
slice access and tape access, the reference-price wall, the as-of regime
probe, the gas and observation builders against future mutation, the sentinel
values, episode truncation and the backtester's own future mutation. All
thirteen pass on the shipped code.

\subsection{Cost models}

Gas is priced at each block's base fee plus a $3\,\%$ tip surcharge, converted
at the block's reference ETH price, and charged as representative per-action
units rather than per-transaction receipts. Table~\ref{tab:app:gas} lists the
units. Slippage is proportional to the notional actually swapped: the pool fee
tier plus a fixed five-basis-point impact. The hedge pays a five-basis-point
taker fee on notional changes and the recorded eight-hour funding rate on the
short.

\begin{table}[tbp]
\caption{Cost-model parameters. Gas units are per action and sweepable in the
configuration; the tip and the slippage impact are the shipped values.}
\label{tab:app:gas}
\centering
\small
\begin{tabular}{@{}L{0.34\linewidth}L{0.22\linewidth}L{0.36\linewidth}@{}}
\toprule
\textbf{Parameter} & \textbf{Value} & \textbf{Charged on} \\
\midrule
Mint gas units & $460\,000$ & opening or rebalancing a position \\
Burn gas units & $215\,000$ & closing the position \\
Collect gas units & $130\,000$ & fee collection \\
Rebalance swap units & $150\,000$ & the swap inside a rebalance \\
Tip surcharge & $3\,\%$ of the base fee & every action \\
Slippage impact & $5$ basis points & the swapped inventory \\
Hedge taker fee & $5$ basis points & short notional adjustments \\
\bottomrule
\end{tabular}
\end{table}

\subsection{The learning interface}

The observation is a fixed-order vector of twenty-four scaled features. Table
\ref{tab:app:observation} lists them in order. The heavy-tailed dimensions are
log-compressed so every input reaches the policy at order one. The regime slot
is the ordinal of the as-of label; the ablation arm in
Appendix~\ref{app:ablations} freezes that slot to a constant.

\begin{table}[tbp]
\caption{The observation vector, in order. The right column gives the feature
in words; the regime slot is the only categorical entry.}
\label{tab:app:observation}
\centering
\small
\begin{tabular}{@{}rrL{0.62\linewidth}@{}}
\toprule
\textbf{\#} & \textbf{Slot} & \textbf{Feature} \\
\midrule
1 & price & pool price \\
2 & sqrt price & square-root price \\
3 & tick & current tick \\
4 & $r_{1h}$ & one-hour reference return \\
5 & $r_{24h}$ & twenty-four-hour reference return \\
6 & $\sigma_{24h}$ & twenty-four-hour realised volatility \\
7 & in range & whether the live band covers the price \\
8 & value & position value in USDC \\
9 & fees0 & uncollected fees in USDC \\
10 & fees1 & uncollected fees in WETH \\
11 & IL & IL against the opening basket, as a fraction of capital \\
12 & liquidity & active pool liquidity \\
13 & fee tier & pool fee tier in pips \\
14 & gas & recent base fee plus priority fee \\
15 & ETH price & recent reference price \\
16 & regime & as-of regime ordinal \\
17 & step & decision index in the episode \\
18 & remaining & decisions left in the episode \\
19 & $d_\ell$ & distance to the lower bound, in tick spacings \\
20 & $d_u$ & distance to the upper bound, in tick spacings \\
21 & width & band width, in tick spacings \\
22 & hedge ratio & current short ratio \\
23 & hedge equity & hedge account equity, as a fraction of capital \\
24 & funding & recent eight-hour funding rate \\
\bottomrule
\end{tabular}
\end{table}

Actions are factored. The general grid pairs nine centre offsets with ten
widths and a hold option; the study's configuration, after the timing finding,
restricts the centre to zero and the width to three arms, and adds the hedge
ratio as an independent head. An action is therefore one of four LP choices
(hold or one of three widths) crossed with five hedge ratios, a
$\mathrm{MultiDiscrete}([4,5])$ space. Table~\ref{tab:app:action} lists the
arms.

\begin{table}[tbp]
\caption{The restricted action grid of the study. Widths are half-widths in
tick spacings; the hedge ratio is the short notional as a fraction of the LP's
WETH delta, and it is chosen at every decision.}
\label{tab:app:action}
\centering
\small
\begin{tabular}{@{}L{0.34\linewidth}L{0.32\linewidth}L{0.28\linewidth}@{}}
\toprule
\textbf{Head} & \textbf{Options} & \textbf{Meaning} \\
\midrule
LP action & hold, deploy $50$, $250$, $1000$ & deploy once, then hold \\
Hedge ratio & $0$, $0.25$, $0.5$, $0.75$, $1.0$ & short share of the delta \\
\bottomrule
\end{tabular}
\end{table}

\subsection{Training settings}

Table~\ref{tab:app:training} lists the optimization settings. Five seeds are
trained per configuration and the checkpoint with the highest validation
Sharpe on the last 180 days of the training lane is the one evaluated. The
configuration was fixed before the evaluation year was scored.

\begin{table}[tbp]
\caption{Optimization settings for every learned policy in the study.}
\label{tab:app:training}
\centering
\small
\begin{tabular}{@{}L{0.42\linewidth}L{0.50\linewidth}@{}}
\toprule
\textbf{Setting} & \textbf{Value} \\
\midrule
Algorithm & clipped PPO with GAE \\
Seeds & $0$ to $4$, trained independently \\
Timesteps per seed & $1\,000\,000$ \\
Discount factor $\gamma$ & $0.9999$ \\
GAE parameter $\lambda_{\mathrm{GAE}}$ & $0.95$ \\
Clip range $\epsilon$ & $0.15$ \\
Learning rate & $10^{-4}$ \\
Batch size & $512$ \\
Epochs per update & $4$ \\
Entropy coefficient & $0.001$ \\
Target KL & $0.01$ \\
Network & two hidden layers of $256$ units \\
Parallel environments & $4$, one process each \\
Checkpoint interval & $100\,000$ steps \\
Selection metric & validation Sharpe, last $180$ days of training \\
Risk penalty & CVaR at $5\,\%$, weight $\lambda=0.05$, one-day window \\
\bottomrule
\end{tabular}
\end{table}
